{{دیوان
| عنوان = اعجاز ہے کسی کا یا گردش زمانہ
| ماخذ = ویکی ماخذ
| ماخذ_لنک = https://ur.wikisource.org/wiki/%D8%A7%D8%B9%D8%AC%D8%A7%D8%B2_%DB%81%DB%92_%DA%A9%D8%B3%DB%8C_%DA%A9%D8%A7_%DB%8C%D8%A7_%DA%AF%D8%B1%D8%AF%D8%B4_%D8%B2%D9%85%D8%A7%D9%86%DB%81
| مقابلہ = اقبال سائبر لائبریری (اقبال اکادمی پاکستان)
| مقابلہ_لنک = https://www.iqbalcyberlibrary.net
}}
<poem>
اعجاز ہے کسی کا یا گردش زمانہ
ٹوٹا ہے ایشیا میں سحر فرنگیانہ

تعمیرِ آشیاں سے میں نے یہ راز پایا
اہلِ نوا کے حق میں بجلی ہے آشیانہ

یہ بندگی خدائی، وہ بندگی گدائی
یہ بندۂ خدا بن یا بندۂ زمانہ

غافل نہ ہو خودی سے کر اپنی پاسبانی
شاید کسی حرم کا تو بھی ہے آستانہ

اے لآ اِلہٰ کے وارث باقی نہیں ہے تجھ میں
گفتارِ دلبرانہ، کردارِ قاہرانہ

تیری نگاہ سے دل سینوں میں کانپتے تھے
کھویا گیا ہے تیرا جذب قلندرانہ

رازِ حرم سے شاید اقبال با خبر ہے
ہیں اس کی گفتگو کے انداز محرمانہ
</poem>
